\section{Related literature}
\label{sec:rewrite-literature}

\paragraph{AI, growth, and automated research.}
My RSI technology builds on the endogenous-growth tradition, in which
research investment produces nonrival knowledge
\citep{romer1990,aghionhowitt1992,jones1995}.
In \citet{romer1990}, markups on new intermediate goods finance research.
Here, a single developer instead invests in the efficiency of AI services
used in both production and research. Recent work asks how AI affects
the economy when AI can automate production, research, or both.
\citet{trammellkorinek2023} distinguish the automation of ordinary production
from the automation of R\&D. \citet{mullainathanrambachan2025} argue that AI may
also reorganize science by bringing idea generation and theory selection into
algorithmic workflows. Endogenous AI development has direct antecedents:
\citet{lu2021ai} models AI as a nonrival input that can contribute to its own
accumulation, with resources allocated to its development.
\citet{gersbachetal2025} model AI as self-learning capital that improves through
application, with knowledge spillovers and skilled labor allocated between AI
and applied research.

\citet{davidsonetal2026} provide a closely related feedback architecture:
better AI makes further AI research more productive, while higher output
expands the resources available for research. Their aggregate-economy
application holds the saving rate and resource-allocation shares fixed.
GATE \citep{erdiletal2025gate}
instead optimizes investment and the allocation of compute between training
and inference through a social planner's problem, with adjustment costs for
compute and conventional capital. These distinctions matter:
endogenous technological progress need not imply a private, profit-maximizing
choice of research spending. In my model, the developer chooses research
compute by comparing its current cost with the value of future monopoly
profit. This decision determines the pace of AI improvement jointly with
household saving, capital accumulation, inference demand, and equilibrium
prices. Research productivity remains a technological assumption; the
resources committed to exploiting it are an intertemporal economic choice.

This mechanism also differs from imposing a path of AI progress.
In their baseline model, \citet{korineksuh2024} study how an exogenous
automation frontier interacts with capital accumulation during a possible
transition to AGI.
\citet{korineketal2026scenarios} map alternative paths of AI capabilities and
adoption into GDP, wages, factor shares, and occupational outcomes through
2030, embedding recursive improvement in the assumed path of AI task
productivity.

\paragraph{Substitution, bottlenecks, and long-run growth.}
Following \citet{davidsonetal2026}, I call an input or task a bottleneck when
its slow improvement constrains aggregate progress even if other inputs or
tasks improve rapidly. \citet{aghionjonesjones2019} apply Baumol's cost-disease
logic to AI and growth: if some essential production tasks remain difficult to
automate, progress in automatable tasks need not translate into comparably rapid
aggregate growth. \citet{jones2025rd} develops the corresponding mechanism
inside R\&D. The effect of AI on research depends jointly on the share of
research tasks AI can perform, its productivity in those tasks, and the
strength of complementarities across tasks. Tasks that AI cannot perform can
therefore limit overall research progress. In my model, human labor no longer enters AI research, but it may remain essential
in final-good production.

The production side connects the CES growth tradition \citep{pitchford1960}
to growth sustained by capital accumulation. \citet{rebelo1991} shows how
reproducible inputs can sustain growth when their production does not require
nonreproducible factors. \citet{jonesmanuelli1990} allow fixed factors to remain
in production, provided the limiting marginal product of capital is high
enough. My AI-dominated regime has this asymptotically linear structure:
capital and AI services can expand together even as AI efficiency approaches
a finite limit. Research determines how efficiency improves; continued
accumulation sustains growth once efficiency is sufficiently high.
\citet{duffypapageorgiou2000} emphasize that a high elasticity is not by itself
sufficient: the productivity of the accumulable input must also be high enough.
In the AI literature, \citet{nordhaus2021} identifies substitutability
between information and conventional inputs as central to an economic
singularity. \citet{jonestonetti2026} endogenize automation in goods and idea
production, treating AI in one scenario as a continuation of historical
automation. Complementary tasks create weak links that delay growth
acceleration. \citet{jones2026future} synthesizes this mechanism, explaining
why initially moderate growth need not rule out a later acceleration.
Complementarity across tasks in that framework is distinct from the elasticity
between labor and AI services in my aggregate production function; the two
parameters are not interchangeable. Consistent with the bottleneck mechanism,
\citet{demirermusolffyang2026} find that large AI-induced
gains in coding activity attenuate sharply when measured by software releases,
and estimate strong complementarity between upstream output and downstream
human effort. Finally, \citet{restrepo2025agi} distinguishes work that remains a
bottleneck from work
that is no longer essential for growth. I obtain a related distinction without
assuming that all work can be performed by AI. Gross substitution between AI
services and labor is necessary but not sufficient for labor to cease being a
bottleneck: attainable AI efficiency must also be high enough. I derive the
threshold separating the labor-bottleneck and AI-dominated equilibria. The
capability threshold in \citet{gans2025aiknowledge} concerns a different
mechanism---whether AI directs research toward dense or exploratory regions of
the knowledge space---rather than whether effective AI services can replace
labor in final production.

\paragraph{Adoption and transition dynamics.}
Technological availability need not translate immediately into measured
productivity gains. \citet{brynjolfssonrocksyverson2021} study complementary
intangible investment and show how its imperfect measurement generates a
productivity J-curve: measured productivity growth initially understates the
gains and later overstates them. Information also matters.
\citet{beaudryportier2014} review how news about future technology changes
current economic decisions. These mechanisms distinguish technological
progress, its implementation, and the response to learning about it. I do
not model organizational capital or compute-adjustment costs. My RSI-activation
exercises isolate an unanticipated change in research opportunities, allowing
consumption and research spending to adjust immediately.

\paragraph{Wages and the labor income share.}
The distributional results connect to the task approach to automation.
\citet{acemoglurestrepo2019} distinguish the displacement effect of automating
existing tasks from two forces that can increase labor demand: the expansion of
output generated by higher productivity and the creation of new tasks in which
labor retains a comparative advantage, which they call the reinstatement effect.
My model does not explicitly model the automation or creation of tasks.
It instead studies how substitution between labor and AI services interacts
with productivity gains.
Labor and AI services enter a constant-elasticity-of-substitution (CES)
production function. Holding capital and effective labor fixed, greater use of
AI services raises labor's income share when the elasticity is below one,
leaves it unchanged when the elasticity equals one, and reduces it when the
elasticity exceeds one. A related CES mechanism underlies
\citet{karabarbounisneiman2014}, who study how a decline in the relative price
of investment induces substitution from labor toward capital and lowers
labor's income share.
Recent work shows that automation can raise wages while reducing labor's
income share \citep{autorkausik2026}. In the long run,
\citet{raymookherjee2022} show that capital accumulation can drive labor's share
to zero while real wages rise indefinitely, and \citet{restrepo2025agi} obtains
a vanishing labor share when compute can perform all bottleneck work.
\citet[p.~14]{jones2026future} also explicitly distinguishes a vanishing
labor income share from falling wages.
The coexistence of rising real wages and a declining labor share is therefore
not itself new. In \citet{farboodikohxia2026}, data generated by economic
activity improve automated tasks and extend automation to new tasks. With
endogenous capital accumulation, this mechanism can generate explosive growth
but stagnant long-run wages. My contribution is a sharper equilibrium
characterization. Among the long-run regimes characterized here with a finite
upper bound on AI efficiency, real-wage growth exceeds the exogenous rate only
in the AI-dominated regime, in which labor's income share converges to
zero. In every positive-labor-share regime, real-wage growth equals the
exogenous rate. Within the AI-dominated regime, an exact expression links the
wage-growth premium to the rate at which labor's income share declines, while
the AI-efficiency threshold determines which regime arises. I do not study
taxation, but a vanishing labor income share connects to the erosion of the
labor-income tax base relative to output analyzed by \citet{korineklockwood2026}.
Their additional concern about a shrinking human-consumption tax base is not
implied by my labor-share result.

The aggregation of labor also matters for interpreting these results.
\citet{korineketal2026scenarios} show that average wages can rise while wages
for cognitive workers fall and their unemployment increases. My
representative-worker model cannot characterize this heterogeneity within
labor. It instead isolates the conditions under which aggregate real wages
rise faster than labor-augmenting technological progress while labor's income
share converges to zero.

\paragraph{Interest rates.}
AI-driven growth can also affect equilibrium returns. Faster expected
consumption growth raises long-term real rates through consumption smoothing in
\citet{chowhalperinmazlish2026}. The AI scenario in \citet{rachel2025} can also
raise the natural rate through saving and investment. In the ICT-investment
scenarios of \citet{melinavilla2025}, the short-run natural rate rises more
when ICT capital complements labor; greater substitutability can weaken or
reverse that response. Using recent AI-sector investment,
\citet{wachterwachter2026} infer scenarios with rapid output and AI-sector
expansion and a higher risk-free rate. I study the long-run net return on
physical capital in a closed-economy Ramsey equilibrium, rather than the
short-run natural rate. The return remains tied to labor-augmenting technological progress
when labor is the bottleneck, but rises with attainable AI efficiency in the
AI-dominated regime. %This is a regime-specific comparative static, not a forecast of market interest rates.

\paragraph{Measurement of AI services.}
The distinction between raw compute and the effective services it produces is
informed by \citet{korinekmckelvey2026}, who distinguish physical computing
capacity from quality-adjusted inference and R\&D/training output and account
for algorithmic progress. I retain separate compute flows for inference and
research but summarize the technology that augments both in one AI-efficiency
stock. \citet{coylepoquiz2025} document the broader difficulties of measuring
AI investment and output, quality change, and changes in production processes.
These measurement problems caution against interpreting the model's
AI-efficiency stock as a directly observed statistical series.

\paragraph{Market structure.}
Private incentives to automate also have important precedents.
\citet{acemoglurestrepo2018} endogenize the allocation of a fixed supply of
scientists between automation and new tasks through expected innovation
profits. \citet{hemousolsen2022} model profit-driven automation and the creation
of new products, with rising low-skill wages strengthening incentives to automate.
Thus, private research incentives are not themselves new. Here they govern
improvements in the efficiency of AI that is used in both production and
further research.

The integrated developer is a tractable limiting case of the concentration
forces studied by \citet{korinekvipra2025} and the AI market-power questions
surveyed by \citet{gans2024}. \citet{atheyscottmorton2025} model AI as a priced
intermediate input and show how upstream market power affects adoption,
downstream prices, factor returns, and welfare. \citet{liuwan2026} compare
decentralized, planner, and monopoly allocations in an endogenous-growth model
linking general-purpose AI to specialized applications. They likewise connect
AI industry structure to innovation incentives. Evidence from API markets on
model pricing, entry, and differentiation in \citet{demireretal2025} also cautions against
treating current AI markets as a literal monopoly. I hold market structure fixed
to isolate the dynamic mechanism. This accounting distinguishes the AI
industry's gross revenue from monopoly profit: inference and research
expenditure are paid before the developer distributes the residual to its
owner.

\paragraph{Contribution.}
The individual mechanisms studied here---automated research, production
bottlenecks, declining labor shares, rising real wages, and changes in
equilibrium returns---appear in the existing literature. I combine them in a
decentralized endogenous-growth model with recursive AI improvement and a
profit-maximizing AI developer. The pace of AI improvement follows from the
developer's research decision, financed against future monopoly profit, rather
than an imposed technological path or a fixed research share. This connects
the transition in wages and income shares to the incentives and resources
supporting AI development. Across the long-run equilibria characterized
here with a finite upper bound on AI efficiency, real wages grow faster than exogenous labor-augmenting technological
progress if and only if labor's income share converges to zero. Output per
person then grows even faster, and an exact expression links the wage-growth
premium to the rate at which labor's income share declines. Within the
AI-dominated regime, higher attainable AI efficiency raises both long-run growth
and the net return on capital.
